\begin{center}
    {\LARGE \textbf{2008分析化学}} \\[2ex]
    {\large 时量180分钟 \quad 满分90分}
\end{center}

\vspace{0.5cm}
\noindent\textbf{注意事项:}
\begin{enumerate}[label=\arabic*., parsep=0pt, itemsep=0pt]
    \item 答题前,考生先将自己的姓名、学号、考场号、座位号填写在答题卡上,将条形码准确粘贴在考生信息条形码粘贴区。
    \item 考试结束后,将本试卷与答题纸一并收回。
    \item 回答各个问题时,将答案写在答题纸上,写在本试卷上无效。
\end{enumerate}

\vspace{0.5cm}
\noindent\textbf{一、简答题(28 pts. total)}

\begin{enumerate}[label=\arabic*., itemsep=1.5ex]
    \item 在含有EDTA的中性溶液中，$\ce{BaSO4}$沉淀的溶解度比在纯水中有所增大，说明原因。
    \item 分析纯的$\ce{SnCl2\cdot H2O}$和$\ce{CaO}$能否作滴定分析的基准物质，为什么？
    \item (1) 将$20\ \mathrm{mL}\ 0.1\ \mathrm{mol/L}\ \ce{NaOH}$溶液和$10\ \mathrm{mL}\ 0.1\ \mathrm{mol/L}\ \ce{H2SO4}$混合，写出所得溶液的质子条件式。
    \\
    (2) 将$20\ \mathrm{mL}\ 0.1\ \mathrm{mol/L}\ \ce{HCl}$和$20\ \mathrm{mL}\ 0.05\ \mathrm{mol/L}\ \ce{Na2CO3}$混合，写出所得溶液的质子条件式。
    \item 如果配制EDTA的水中含有$\ce{Ca^{2+}}$，$\ce{Mg^{2+}}$，对下列结果会产生什么影响（偏高、偏低、无影响）？
    \begin{enumerate}[label=(\arabic*), noitemsep, topsep=0pt, partopsep=0pt]
        \item 在碱性溶液中标定EDTA的浓度
        \item 在碱性溶液中测定其它金属离子的浓度
    \end{enumerate}
    \item 碘量法的主要误差来源是什么？
    \item 指出$\ce{Fe(OH)3}$的沉淀类型，洗涤该沉淀应用什么资源？
    \item 分析化学中常用的分离方法有哪些？
\end{enumerate}

\vspace{0.5cm}
\noindent\textbf{二、计算题(62 pts. total)}

\begin{enumerate}[label=\arabic*., itemsep=1.5ex]
    \item (10 pts.) 某试剂水溶液$50\ \mathrm{mL}$，若将$99\%$的有效成分萃取到苯中，问：(1) 用等体积苯萃取一次，分配比$D$需多少才行？(2) 每次用$25.0\ \mathrm{mL}$苯萃取二次，分配比$D$又为多少才行？
    \item (10 pts.) 已知氧化还原反应
    \[
    \ce{MnO^-_{4} + 5Fe^{2+} + 8H^{+} = Mn^{2+} + 5Fe^{3+} + 4H_{2}O}
    \]
    在$1\ \mathrm{mol/L}\ \ce{HClO4}$介质中两电对的电极电位分别为$1.45\ \mathrm{V}$，$0.75\ \mathrm{V}$。计算(1)计量点时的电位 (2)计量点前后$0.1\%$时体系的电位。
    \item (12 pts.) 用$0.20\ \mathrm{mol/L}\ \ce{NaOH}$滴定$0.20\ \mathrm{mol/L}\ \ce{HCl}$（其中含有$0.10\ \mathrm{mol/L}\ \ce{NH4Cl}$）。（已知$\mathrm{p}K_\mathrm{a}(\ce{NH4^+})=9.26$）
    \begin{enumerate}[label=(\arabic*), noitemsep, topsep=0pt, partopsep=0pt]
        \item 计算化学计量点的$\mathrm{pH}$
        \item 若滴定至$\mathrm{pH}=7.0$，问终点时有百分之几的$\ce{NH4^+}$变为$\ce{NH3}$；
        \item 此时终点误差是多少？
    \end{enumerate}
    \item (10 pts.) 于$\mathrm{pH}=5.5$时，用$0.020\ \mathrm{mol/L}\ \mathrm{EDTA}$滴定$0.020\ \mathrm{mol/L}\ \ce{Zn^{2+}}$和$0.20\ \mathrm{mol/L}\ \ce{Mg^{2+}}$的混合溶液中的$\ce{Zn^{2+}}$，(1) 计算化学计量点时$\ce{Zn^{2+}}$和$\ce{MgY^{2-}}$的浓度；(2) 计算终点误差（以二甲酚橙为指示剂）。
    \\
    ［已知$\mathrm{pH}=5.5$时，$\lg\alpha_{\mathrm{Y(H)}}=5.5$，$\mathrm{pZn_{(二甲酚橙)}}=5.7$，$K(\ce{ZnY})=10^{16.5}$，$K(\ce{MgY})=10^{8.70}$］
    \item (10 pts.) 一组测量值为：$20.04\%$，$20.01\%$，$20.04\%$，$20.07\%$，$20.08\%$。若真值为$20.10\%$，计算$95\%$置信度的平均值置信区间。此区间是否包括真值在内？若置信度定为$99\%$，是否包括真值在内？计算结果说明什么问题？
    \\
    ［已知$t_{0.05,4}=2.78$，$t_{0.05,5}=2.57$，$t_{0.01,4}=4.60$，$t_{0.01,5}=4.03$］
    \item (10 pts.) 计算$\ce{AgCl}$沉淀在$\mathrm{pH}=8.00$，络合剂$\ce{L}$的总浓度$c(\ce{L})=0.10\ \mathrm{mol/L}$溶液中的溶解度。（忽略在形成络合物$\ce{L}$的消耗）
    \\
    ［已知$K_\mathrm{a}(\ce{HL})=1.0\times10^{-10}$，$\ce{AgL2}$的$\lg\beta_1=3.0$，$\lg\beta_2=7.0$，$K_{\mathrm{sp}}(\ce{AgCl})=1.8\times10^{-10}$］
\end{enumerate}